\documentclass[12pt]{article}
\usepackage[utf8]{inputenc}
\usepackage[T1]{fontenc}
\usepackage[margin=1in]{geometry}
\usepackage{setspace}
\usepackage{xcolor}
\usepackage[colorlinks=true, urlcolor=blue!60!black, linkcolor=blue!60!black]{hyperref}
\setlength{\parindent}{2em}
\setlength{\parskip}{1.25em}
\renewcommand{\baselinestretch}{1.0}
\begin{document}
\begin{center}
{\LARGE \textit{Everson v. Board of Education}, 330 U.S. 1 (1947)\par}
\vspace{0.8em}
{\large Case Brief\par}
\end{center}
\vspace{1em}

\section{Memory Jogger}
Incorporating the Establishment Clause against the states and invoking Jefferson's ``wall of separation between church and State,'' the Court nonetheless upheld 5--4 New Jersey's reimbursement of bus fares for children attending parochial schools as a neutral public-welfare benefit, like police and fire protection (\href{https://supreme.justia.com/cases/federal/us/330/1/}{Justia}).

\section{Detailed Case Facts}
In 1941, the New Jersey Legislature enacted a statute (N.J. Laws 1941, c. 191, p. 581; N.J. Rev. Stat. Cum. Supp. tit. 18, c. 14, \S{} 8) authorizing local school district boards of education to make rules and contracts for transporting children living remote from any schoolhouse to and from school, ``including the transportation of school children to and from school other than a public school, except such school as is operated for profit in whole or in part.'' Acting under this statutory authority, the Board of Education of the Township of Ewing---a suburban township near Trenton that operated public elementary schools only through the eighth grade and maintained no public high school---adopted a resolution authorizing the semiannual reimbursement of parents from general tax-raised school funds for fares paid on regular buses operated by the public transportation system.

Under the Ewing Township resolution, reimbursement was expressly authorized for pupils traveling by public carrier to three public high schools outside the district (Trenton Senior High School, Trenton Junior High School, and Pennington High School) and four Catholic parochial schools (St. Mary's Cathedral High School, Trenton Catholic Boys High School, St. Hedwig's Parochial School, and St. Francis School). As established in the record by the superintendent of parochial schools for the Diocese of Trenton---a Catholic priest---these parochial schools satisfied New Jersey's secular educational requirements while also providing regular religious instruction conforming to the tenets and worship of the Roman Catholic Church. Reimbursements averaged roughly \$40 per child annually.

\section{Procedural History}
Arch R. Everson, a taxpayer residing in Ewing Township, filed suit in New Jersey state court challenging the constitutionality of both the 1941 statute and the township board's resolution insofar as they authorized the expenditure of tax-raised funds to reimburse parents for transporting children to parochial schools. The New Jersey Supreme Court (then an intermediate appellate court) held that the legislature lacked authority under the New Jersey Constitution to authorize reimbursement of bus fares for nonpublic school students, 132 N.J.L. 98, 39 A.2d 75 (1944). On appeal, the New Jersey Court of Errors and Appeals (then the state's highest court) reversed, holding that neither the statute nor the resolution violated the New Jersey Constitution or the Federal Constitution, 133 N.J.L. 350, 44 A.2d 333 (1945). Everson appealed the federal constitutional questions to the United States Supreme Court under 28 U.S.C. \S{} 344(a). Notably, Everson did not press an Equal Protection Clause challenge to the statute's exclusion of for-profit private schools, nor did the record establish that any Ewing Township children attended non-Catholic private or religious schools.

\section{Judicial Votes}
The Supreme Court \textbf{affirmed} the New Jersey Court of Errors and Appeals by a \textbf{5--4} vote. The case was argued on November 20, 1946, and decided on February 10, 1947 (\href{https://supreme.justia.com/cases/federal/us/330/1/}{Justia}; \href{https://www.law.cornell.edu/supremecourt/text/330/1}{Cornell LII}):

\begin{itemize}
    \item \textbf{Majority Opinion:} Justice Black delivered the opinion of the Court, joined by Chief Justice Vinson and Justices Reed, Douglas, and Murphy.
    \item \textbf{Concurring Opinions:} None.
    \item \textbf{Dissenting Opinions:} Justice Jackson filed a dissenting opinion, joined by Justice Frankfurter. Justice Rutledge filed a dissenting opinion, joined by Justices Frankfurter, Jackson, and Burton.
\end{itemize}

Notably, all four dissenting Justices agreed with Justice Black's two foundational constitutional premises---that the Establishment Clause is incorporated against the states via the Fourteenth Amendment and that it mandates a strict separation of church and state---but forcefully argued that the majority betrayed its own principles in upholding the Ewing Township reimbursement scheme.

\section{Holding}
\textbf{Questions Presented:}

\begin{enumerate}
  \item \textbf{Fourteenth Amendment Due Process (Public Purpose):} Does a state statute and local school board resolution authorizing the reimbursement of parents from tax-raised funds for public bus fares paid to transport their children to parochial schools constitute an unconstitutional taking of private property by taxation for a private, rather than a public, purpose in violation of the Due Process Clause of the Fourteenth Amendment?
  \item \textbf{Incorporation of the Establishment Clause:} Does the Fourteenth Amendment incorporate the First Amendment's Establishment Clause so as to make its prohibition against ``any law respecting an establishment of religion'' binding upon the states?
  \item \textbf{Establishment Clause Merits:} Assuming the Establishment Clause applies to the states, does New Jersey's general program reimbursing parents for public bus transportation to accredited public and nonprofit parochial schools violate the First Amendment's prohibition against an establishment of religion?
\end{enumerate}

\textbf{Holding:} In a 5--4 decision authored by Justice Hugo L. Black, the Supreme Court \textbf{affirmed} the judgment of the New Jersey Court of Errors and Appeals and held:

\begin{enumerate}
  \item \textbf{No Due Process Violation:} Using tax-raised funds to reimburse parents for children's bus fares to accredited schools---including parochial schools---serves a legitimate public purpose (facilitating secular education and protecting children from traffic hazards) and therefore does not violate the Due Process Clause of the Fourteenth Amendment.
  \item \textbf{Incorporation of the Establishment Clause:} The First Amendment's Establishment Clause is incorporated against the states through the Fourteenth Amendment and imposes the same broad constitutional restrictions on state governments as on the Federal Government.
  \item \textbf{No Establishment Clause Violation:} Although the Establishment Clause erects a ``high and impregnable'' ``wall of separation between church and State'' prohibiting government from aiding one religion, aiding all religions, or levying any tax to support religious institutions, New Jersey did not breach that wall. By extending a neutral, generally available public-welfare benefit to parents regardless of their faith---comparable to police and fire protection---the state remained neutral rather than adversarial toward religion and contributed no money directly to the parochial schools.
\end{enumerate}

\section{Analysis of Opinions}
\subsection{Majority Opinion (Justice Black)}

\subsection{Threshold Framing and the Fourteenth Amendment Public-Purpose Inquiry}

Justice Black first set aside the statute's exclusion of for-profit private schools because the appellant had raised no Equal Protection challenge and the record contained no evidence that any Ewing Township children attended, or would have attended, any school other than public or Catholic schools (\href{https://supreme.justia.com/cases/federal/us/330/1/}{330 U.S. at 4--5 \& n.2}). Turning to the Due Process ``public purpose'' objection, the Court reaffirmed that federal judicial authority to invalidate state tax expenditures for lacking a public purpose must be exercised with ``the most extreme caution'' so as not to curtail states' core police power to legislate for the general welfare (\href{https://supreme.justia.com/cases/federal/us/330/1/}{330 U.S. at 6--7}). Relying on \href{https://supreme.justia.com/cases/federal/us/281/370/}{\emph{Cochran v. Louisiana State Board of Education}, 281 U.S. 370 (1930)} (which had upheld state provision of secular textbooks to private and parochial school students on a ``child-benefit'' public-purpose rationale under the Due Process Clause), Justice Black concluded that it was ``much too late to argue that legislation intended to facilitate the opportunity of children to get a secular education serves no public purpose'' (\href{https://supreme.justia.com/cases/federal/us/330/1/}{330 U.S. at 7}). Reimbursing parents so children could ride public buses rather than risk traffic hazards walking or ``hitchhiking'' remained a public program even though it coincided with parents' private desires.

\subsection{Incorporation and the Historical Scope of the Establishment Clause}

Addressing Everson's central claim, the Court squarely incorporated the Establishment Clause against the states via the Fourteenth Amendment. Citing \href{https://supreme.justia.com/cases/federal/us/319/105/}{\emph{Murdock v. Pennsylvania}, 319 U.S. 105 (1943)}, and the Court's Free Exercise incorporation precedents (including \emph{Cantwell v. Connecticut}, 310 U.S. 296 (1940)), Justice Black stated without extended doctrinal debate that there was ``every reason to give the same application and broad interpretation to the `establishment of religion' clause'' (\href{https://supreme.justia.com/cases/federal/us/330/1/}{330 U.S. at 8, 15}).

To define the clause's substantive meaning, Justice Black turned to European and colonial history, emphasizing religious persecution, compulsory church attendance, and forced tithes in the American colonies. He treated the 1785--1786 Virginia disestablishment struggle---specifically James Madison's \emph{Memorial and Remonstrance Against Religious Assessments} and Thomas Jefferson's \emph{Virginia Bill for Religious Liberty}---as the decisive interpretive lens for the First Amendment (\href{https://supreme.justia.com/cases/federal/us/330/1/}{330 U.S. at 11--13}). From this history, Black distilled his famous list of strict-separationist prohibitions: neither a state nor the federal government may set up a church, pass laws that aid one religion or all religions, coerce belief or attendance, punish religious belief or disbelief, levy ``[n]o tax in any amount, large or small'' to support religious activities or institutions, or participate in the affairs of religious organizations (\href{https://supreme.justia.com/cases/federal/us/330/1/}{330 U.S. at 15--16}). Invoking Jefferson's 1802 letter to the Danbury Baptists (previously quoted in \emph{Reynolds v. United States}, 98 U.S. 145, 164 (1879)), Black declared that the clause was intended to erect ``a wall of separation between church and State'' (\href{https://supreme.justia.com/cases/federal/us/330/1/}{330 U.S. at 16}).

\subsection{Reconciling Strict Separation with Neutral Public-Welfare Benefits}

Having built a seemingly absolute separationist framework, Justice Black pivoted sharply to uphold the New Jersey program, noting that the Court must not strike down a state statute ``even though it approaches the verge of [the State's constitutional] power'' (\href{https://supreme.justia.com/cases/federal/us/330/1/}{330 U.S. at 16}). Black identified a structural tension between the two Religion Clauses: while the Establishment Clause forbids contributing tax funds to support religious institutions, the Free Exercise Clause forbids hampering citizens' religious liberty. Consequently, a state ``cannot exclude individual Catholics, Lutherans, Mohammedans, Baptists, Jews, Methodists, Nonbelievers, Presbyterians, or the members of any other faith, because of their faith, or lack of it, from receiving the benefits of public welfare legislation'' (\href{https://supreme.justia.com/cases/federal/us/330/1/}{330 U.S. at 16}).

Black candidly acknowledged that bus reimbursements helped children reach church schools and might influence some parents' enrollment choices. Yet he analogized the program to general public safety and infrastructure services---such as traffic policemen stationed near parochial schools, ordinary police and fire protection, sewage connections, and public highways and sidewalks---that are ``so separate and so indisputably marked off from the religious function'' that withholding them would make the state an adversary of religion rather than a neutral actor (\href{https://supreme.justia.com/cases/federal/us/330/1/}{330 U.S. at 17--18}). Invoking \href{https://supreme.justia.com/cases/federal/us/268/510/}{\emph{Pierce v. Society of Sisters}, 268 U.S. 510 (1925)}, which protected parents' constitutional right to satisfy compulsory education laws at accredited religious schools, Black concluded that New Jersey contributed no money to the schools themselves, but merely operated a neutral safety program helping parents get their children safely to accredited schools.

\subsection{Dissenting Opinion (Justice Jackson)}

Justice Jackson opened with a celebrated indictment of the internal contradiction between Black's strict-separationist rhetoric and his permissive holding, comparing the Court to Byron's Julia who, ``whispering `I will ne'er consent,' -- consented'' (\href{https://supreme.justia.com/cases/federal/us/330/1/}{330 U.S. at 19}). Jackson attacked the majority on three principal grounds:

\begin{itemize}
  \item \textbf{Record Distortion and Denominational Discrimination:} Ewing Township did not operate school buses or furnish a safety service; children rode ordinary public buses ``as fast, and no faster, and \ldots{} as safe, and no safer'' after reimbursement (\href{https://supreme.justia.com/cases/federal/us/330/1/}{330 U.S. at 20}). Worse, the township resolution reimbursed only public and Catholic school students while the statute excluded for-profit private schools. Unlike a police officer or firefighter who protects people and buildings without asking about religious affiliation, Ewing officials had to ask a ``religious test'' question---``is this \ldots{} school identified with the Catholic Church?''---before issuing a check (\href{https://supreme.justia.com/cases/federal/us/330/1/}{330 U.S. at 25--26}).
  \item \textbf{Inseparability of Parochial Schools from the Church:} Quoting Catholic Canon Law (Canons 1372--1374 and 1381), Jackson argued that parochial education is ``the rock on which the whole structure [of the Roman Catholic Church] rests,'' making tax aid to parochial schools indistinguishable from tax aid to the Church itself (\href{https://supreme.justia.com/cases/federal/us/330/1/}{330 U.S. at 22--24}).
  \item \textbf{Regulatory Entanglement Warning:} Jackson warned religious schools that public subsidies inevitably invite state control: ``If the state may aid these religious schools, it may therefore regulate them,'' quoting \emph{Wickard v. Filburn}, 317 U.S. 111, 131 (1942), for the proposition that ``It is hardly lack of due process for the Government to regulate that which it subsidizes'' (\href{https://supreme.justia.com/cases/federal/us/330/1/}{330 U.S. at 27--28}).
\end{itemize}

\subsection{Dissenting Opinion (Justice Rutledge)}

Justice Rutledge authored an exhaustive, historically anchored dissent---appending the complete texts of James Madison's 1785 \emph{Memorial and Remonstrance Against Religious Assessments} and the defeated 1784 Virginia Assessment Bill (\href{https://supreme.justia.com/cases/federal/us/330/1/}{330 U.S. at 63--74}). Rutledge argued:

\begin{itemize}
  \item \textbf{Unitary, Comprehensive Scope of ``Religion'':} Because the word ``religion'' appears only once in the First Amendment and governs both clauses via ``thereof,'' the Establishment Clause is just as broad as the Free Exercise Clause; it prohibits not merely a single established church, but ``every form of public aid or support for religion'' (\href{https://supreme.justia.com/cases/federal/us/330/1/}{330 U.S. at 31--33}).
  \item \textbf{Madison's ``Three Pence'' Principle and Non-Discriminatory Aid:} Rutledge demonstrated that the 1784 Virginia Assessment Bill had been modified to allow taxpayers to designate any Christian church---or even general education---for their tax payments, yet Madison fiercely opposed even non-discriminatory assessments. Multiplying support to all faiths does not cure the constitutional wrong of forcing citizens to finance ``the propagation of opinions which [they] disbelieve'' (\href{https://supreme.justia.com/cases/federal/us/330/1/}{330 U.S. at 36--41, 45--46, 59--60}).
  \item \textbf{Collapse of Any Limiting Principle:} Rejecting the ``child-benefit'' and ``public-welfare'' rationales as an improper transplant from Due Process jurisprudence, Rutledge contended that transportation is just as essential to religious schooling as textbooks, lunches, buildings, tuition, or teachers' salaries. Once transportation is deemed a ``public function,'' ``[n]o rational line can be drawn between payment for such larger, but not more necessary, items and payment for transportation'' (\href{https://supreme.justia.com/cases/federal/us/330/1/}{330 U.S. at 47--53}). Hardship to religious parents who pay public school taxes while funding parochial schools is the intentional ``great price'' of religious liberty (\href{https://supreme.justia.com/cases/federal/us/330/1/}{330 U.S. at 58--59}).
\end{itemize}

\subsection{Subsequent Developments}

\begin{itemize}

\item \textbf{Early Applications of the Wall and the Child-Benefit Theory (1948--1971):} Early decisions enforced Justice Black's ``wall'' to bar in-school religious instruction in \href{https://supreme.justia.com/cases/federal/us/333/203/}{\emph{McCollum v. Board of Education}, 333 U.S. 203 (1948)}, and public school prayer in \href{https://supreme.justia.com/cases/federal/us/370/421/}{\emph{Engel v. Vitale}, 370 U.S. 421 (1962)}, while allowing off-campus released time in \href{https://supreme.justia.com/cases/federal/us/343/306/}{\emph{Zorach v. Clauson}, 343 U.S. 306 (1952)}, and secular textbook loans in \href{https://supreme.justia.com/cases/federal/us/392/236/}{\emph{Board of Education v. Allen}, 392 U.S. 236 (1968)}. In \href{https://supreme.justia.com/cases/federal/us/403/602/}{\emph{Lemon v. Kurtzman}, 403 U.S. 602 (1971)}, the Court formalized its three-prong purpose, effect, and entanglement test to block direct parochial school salary supplements.

\item \textbf{Neutrality Becomes a Mandate (2002--2022):} Since 2002, however, the Court has elevated \emph{Everson}'s neutrality strand and dismantled its ``no-aid'' wall. \href{https://supreme.justia.com/cases/federal/us/536/639/}{\emph{Zelman v. Simmons-Harris}, 536 U.S. 639 (2002)}, upheld neutral school vouchers directed by ``true private choice.'' Although \href{https://supreme.justia.com/cases/federal/us/540/712/}{\emph{Locke v. Davey}, 540 U.S. 712 (2004)}, preserved ``play in the joints'' to bar state scholarships for devotional ministry degrees, the Court's Free Exercise trilogy turned \emph{Everson}'s nondiscrimination dictum into a binding mandate: \href{https://supreme.justia.com/cases/federal/us/582/15-577/}{\emph{Trinity Lutheran Church of Columbia, Inc. v. Comer}, 582 U.S. 449 (2017)} (playground grants), \href{https://supreme.justia.com/cases/federal/us/591/18-1195/}{\emph{Espinoza v. Montana Department of Revenue}, 591 U.S. 464 (2020)} (invalidating Blaine Amendment exclusions from tax-credit scholarships), and \href{https://supreme.justia.com/cases/federal/us/596/20-1088/}{\emph{Carson v. Makin}, 596 U.S. 767 (2022)} (erasing the status/use distinction for Maine tuition aid).

\end{itemize}

\section{Examples: Future Applications}
\begin{enumerate}

\item \textbf{Same-Side (Aid Upheld): Free Transit Passes for All Students.} A city gives every enrolled K--12 student a free pass for its public buses and subways, whether the student attends a public, private secular, or religious school. A taxpayer sues, arguing that passes used to reach Catholic and Jewish day schools are tax support for religious institutions. Under \textit{Everson}, the program is a general public-welfare benefit paid to students and their families, not to the schools, and it is indistinguishable from the bus-fare reimbursements the Court upheld. Excluding religious-school students would make the city an adversary of religion rather than a neutral, so the passes survive the Establishment Clause challenge (\href{https://supreme.justia.com/cases/federal/us/330/1/}{Justia}).

\item \textbf{Same-Side (Aid Upheld): Loaning Secular Textbooks to All Students.} A state lends approved secular textbooks in mathematics, science, and history to every student in grades 7--12, including students at parochial schools; the books remain state property and the school never receives money. Taxpayers object that the books free up parochial-school budgets for religious instruction. Extending \textit{Everson}'s child-benefit reasoning, \href{https://supreme.justia.com/cases/federal/us/392/236/}{\textit{Board of Education v.\ Allen}, 392 U.S.\ 236 (1968)}, upheld exactly this kind of program: the benefit runs to students and parents, the materials are secular, and the State contributes nothing to the schools' religious mission. Rutledge's fungibility objection did not carry the day in \textit{Everson}, and it does not here.

\item \textbf{Opposite-Side (Aid Invalid): Tax Grants to Pay Clergy Salaries.} To help congregations hit hard by an economic downturn, a county appropriates tax revenue for grants to churches, synagogues, and mosques to pay the salaries of their clergy, distributing the money equally among all faiths that apply. Equal treatment of all religions does not save the program. \textit{Everson} declared that government may not pass laws that aid all religions and that no tax in any amount may be levied to support religious activities or institutions, and the grants pay for the core religious function itself rather than a service ``separate'' from it like police or fire protection. Madison's objection to the nondiscriminatory Virginia assessment, which both the majority and Rutledge embraced, applies directly.

\item \textbf{Opposite-Side (Practice Invalid): Religious Classes Inside Public Schools.} A public school district lets clergy of several faiths come into its classrooms for an hour each week during the school day to teach religion to students whose parents consent, while other students move to another room for secular study. The program uses the public school building and the State's compulsory-attendance machinery to help religious groups spread their faith. Applying \textit{Everson}'s incorporated Establishment Clause, \href{https://supreme.justia.com/cases/federal/us/333/203/}{\textit{McCollum v.\ Board of Education}, 333 U.S.\ 203 (1948)}, held such a program unconstitutional; unlike \textit{Everson}'s bus fares, the aid here goes to religious instruction itself.

\item \textbf{Fence-Sitter (Unresolved): Private Aid Versus Charter Status.} The State of Franklin reimburses rural parents for private high school tuition and busing, excluding ``sectarian'' schools (Track A), and contracts with private nonprofits to run tuition-free ``public charter schools'' that must be nonsectarian (Track B). Under \emph{Everson}, \emph{Zelman}, and \emph{Carson}, Track A violates the Free Exercise Clause. Track B turns on the unresolved \emph{Drummond} state-action threshold: if charter schools are state actors, the Establishment Clause bars sectarian operation; if they are private benefit recipients, \emph{Carson} forbids religious exclusion.

\end{enumerate}

\section{Critique}
\textbf{Federalism and History:} Critics contend that the original Establishment Clause barred federal interference with state establishments, making casual incorporation problematic, and fault Black and Rutledge for treating Virginia's 1785--1786 struggle as the sole gloss on 1789.

\textbf{Fungibility and Regulatory Strings:} As Rutledge noted, bus reimbursements free up funds for religious instruction; conversely, Jackson's \emph{Wickard} warning---that government regulates what it subsidizes---now drives post-\emph{Carson} disputes over conditions on public aid.

\section{Key Quotations}
\begin{quote}
``\,The ``establishment of religion'' clause of the First Amendment means at least this: neither a state nor the Federal Government can set up a church. Neither can pass laws which aid one religion, aid all religions, or prefer one religion over another. Neither can force nor influence a person to go to or to remain away from church against his will or force him to profess a belief or disbelief in any religion. No person can be punished for entertaining or professing religious beliefs or disbeliefs, for church attendance or non-attendance. No tax in any amount, large or small, can be levied to support any religious activities or institutions, whatever they may be called, or whatever form they may adopt to teach or practice religion. Neither a state nor the Federal Government can, openly or secretly, participate in the affairs of any religious organizations or groups, and vice versa. In the words of Jefferson, the clause against establishment of religion by law was intended to erect ``a wall of separation between church and State.''\,''
\end{quote}
\noindent --- Black, J., for the Court, \href{https://supreme.justia.com/cases/federal/us/330/1/}{330 U.S. at 15--16}.

\begin{quote}
``That Amendment requires the state to be a neutral in its relations with groups of religious believers and nonbelievers; it does not require the state to be their adversary. State power is no more to be used so as to handicap religions than it is to favor them. \ldots{} The First Amendment has erected a wall between church and state. That wall must be kept high and impregnable. We could not approve the slightest breach. New Jersey has not breached it here.''
\end{quote}
\noindent --- Black, J., for the Court, \href{https://supreme.justia.com/cases/federal/us/330/1/}{330 U.S. at 18}.

\begin{quote}
``In fact, the undertones of the opinion, advocating complete and uncompromising separation of Church from State, seem utterly discordant with its conclusion, yielding support to their commingling in educational matters. The case which irresistibly comes to mind as the most fitting precedent is that of Julia who, according to Byron's reports, ``whispering `I will ne'er consent,' -- consented.''\,''
\end{quote}
\noindent --- Jackson, J., dissenting, \href{https://supreme.justia.com/cases/federal/us/330/1/}{330 U.S. at 19}.

\begin{quote}
``But we cannot have it both ways. Religious teaching cannot be a private affair when the state seeks to impose regulations which infringe on it indirectly, and a public affair when it comes to taxing citizens of one faith to aid another, or those of no faith to aid all. \ldots{} If the state may aid these religious schools, it may therefore regulate them. Many groups have sought aid from tax funds, only to find that it carried political controls with it.''
\end{quote}
\noindent --- Jackson, J., dissenting, \href{https://www.law.cornell.edu/supremecourt/text/330/1}{330 U.S. at 27--28}.

\begin{quote}
``Neither so high nor so impregnable today as yesterday is the wall raised between church and state by Virginia's great statute of religious freedom and the First Amendment, now made applicable to all the states by the Fourteenth. New Jersey's statute sustained is the first, if indeed it is not the second, breach to be made by this Court's action. That a third, and a fourth, and still others will be attempted we may be sure.''
\end{quote}
\noindent --- Rutledge, J., dissenting, \href{https://supreme.justia.com/cases/federal/us/330/1/}{330 U.S. at 29}.


\end{document}
